% Ad Familiares 15.1 (ad-familiares-15-01)
% Source: https://raw.githubusercontent.com/PerseusDL/canonical-latinLit/master/data/phi0474/phi056/phi0474.phi056.perseus-lat1.xml
% Fetched by scripts/fetch_latin.py

% Dateline (Perseus): Scr. in Cilicia xiii K. Oct. a. 703 (51).

\ciceroLetterOpener{M. TVLLIVS M. F. CICERO PROCOS. S. D. COS. PR., TR. PL., SENATVI.}

\ciceroSection{1} S. v. v. h. e. e. q. v. etsi non dubie mihi nuntiabatur Parthos transisse Euphratem cum omnibus fere suis copiis, tamen, quod arbitrabar a M. Bibulo procos. certiora de his rebus ad vos scribi posse, statuebam mihi non necesse esse publice scribere ea quae de alterius provincia nuntiarentur. postea vero quam certissimis auctoribus, legatis, nuntiis, 10 litteris, sum certior factus, vel quod tanta res erat vel quod nondum audieramus Bibulum in Syriam venisse vel quia administratio huius belli mihi cum Bibulo paene est communis, quae ad me delata essent scribenda ad vos putavi.

\ciceroSection{2} Regis Antiochi Commageni legati primi mihi nuntiarunt is Parthorum magnas copias Euphratem transire coepisse. quo nuptio adlato, cum essent non nulli qui ei regi minorem fidem habendam putarent, statui exspectandum esse si quid certius adferretur. A. d. xiii K. Oct. , cum exercitum in Ciliciam ducerem, in finibus Lycaoniae et Cappadociae mihi litterae redditae sunt a Tarcondimoto, qui fidelissimus socius trans Taurum amicissimusque p. R. existimatur, Pacorum Orodi regis Parthorum filium cum permagno equitatu Parthico transisse Euphratem et castra posuisse Tybae, magnumque tumultum esse in provincia Syria excitatum. eodem die ab Iamblicho, phylarcho Arabum, quem homines opinantur bene sentire amicumque esse rei p. nostrae, litterae de isdem rebus mihi redditae sunt.

\ciceroSection{3} his rebus adlatis, etsi intellegebam socios infirme animatos esse et novarum rerum exspectatione suspensos, sperabam tamen eos, ad quos iam accesseram quique nostram mansuetudinem integritatemque perspexerant, amiciores p. R. esse factos, Ciliciam autem firmiorem fore, si aequitatis nostrae particeps facta esset. et ob eam causam et ut opprimerentur ii qui ex Cilicum gente in armis essent, et ut hostis is qui esset in Syria sciret exercitum p. R. non modo non cedere iis nuntiis adlatis sed etiam propius accedere, exercitum ad Taurum institui ducere.

\ciceroSection{4} sed si quid apud vos auctoritas mea ponderis habet, in iis praesertim rebus, quas vos audistis, ego paene cerno, magno opere vos et hortor et moneo ut his provinciis serius vos quidem quam decuit, sed aliquando tamen consulatis. nos quem ad modum instructos et quibus praesidiis munitos ad tanti belli opinionem miseritis non estis ignari. quod ego negotium non stultitia occaecatus sed verecundia deterritus non recusavi ; neque enim umquam ullum periculum tantum putavi quod subterfugere mallem ' quam vestrae auctoritati obtemperare.

\ciceroSection{5} hoc autem tempore res sese sic habet ut, nisi exercitum tantum, quantum ad maximum bellum mittere soletis, mature in has provincias miseritis, summum periculum sit ne amittendae sint omnes hae provinciae, quibus vectigalia p. R. continentur. quam ob rem autem in hoc provinciali dilectu spem habeatis aliquam causa nulla est. neque multi sunt et diffugiunt qui sunt metu oblato ; et quod genus hoc militum sit iudicavit vir fortissimus M. Bibulus in Asia, qui, cum vos ei permisissetis, dilectum habere noluerit. nam sociorum auxilia propter acerbitatem atque iniurias imperi nostri aut ita imbecilla sunt ut non multum nos iuvare possint, aut ita alienata a nobis ut neque exspectandum ab iis neque committendum iis quicquam esse videatur.

\ciceroSection{6} Regis Deiotari et voluntatem et copias, quantaecumque sunt, nostras esse duco ; Cappadocia est inanis, reliqui reges tyrannique neque opibus satis firmi nec voluntate sunt. mihi in hac paucitate ' militum animus certe non deerit, spero ne consilium quidem ; quid casurum sit incertum est. utinam saluti nostrae consulere possimus! dignitati certe consulemus.
